% Part: incompleteness
% Chapter: representability-in-q
% Section: representable-comp

\documentclass[../../../include/open-logic-section]{subfiles}

\begin{document}

\olfileid{inc}{req}{rpc}
\olsection{توابع بازنمایی‌پذیر در~$\Th{Q}$ محاسبه‌پذیرند}

ثابت خواهیم کرد هر تابعی که در~$\Th{Q}$ بازنمایی‌پذیر باشد
محاسبه‌پذیر است. نخست باید لمی دربارهٔ توابع بازنمایی‌پذیر در
$\Th{Q}$ برقرار کنیم.

\begin{lem}\ollabel{lem:rep-q}
  اگر $f(x_0,\dots,x_k)$ در~$\Th{Q}$ بازنمایی‌پذیر باشد،
  !!a{formula}ی چون~$!A(x_0,\dots,x_k,y)$ وجود دارد به‌گونه‌ای که
  \[
  \Th{Q} \Proves !A_f(\num{n_0}, \dots, \num{n_k}, \num{m})
  \quad\text{اگر و تنها اگر}\quad m = f(n_0, \dots, n_k).
  \]
\end{lem}

\begin{proof}
جهت «اگر» همان
\olref[int]{defn:representable-fn}\olref[int]{defn:rep:a} است. جهت
«تنها اگر» را چنین می‌بینیم: فرض کنید
$\Th{Q}\Proves !A_f(\num{n_0},\dots,\num{n_k},\num{m})$، اما
$m\neq f(n_0,\dots,n_k)$. بگذارید
$l=f(n_0,\dots,n_k)$. بنا بر
\olref[int]{defn:representable-fn}\olref[int]{defn:rep:a}،
$\Th{Q}\Proves !A_f(\num{n_0},\dots,\num{n_k},\num{l})$. بنا بر
\olref[int]{defn:representable-fn}\olref[int]{defn:rep:b}،
$\Th{Q}\Proves\lforall[y][(!A_f(\num{n_0},\dots,\num{n_k},y)
\lif\num{l}=y)]$. با استفاده از منطق و فرض
$\Th{Q}\Proves !A_f(\num{n_0},\dots,\num{n_k},\num{m})$، نتیجه
می‌گیریم $\Th{Q}\Proves\eq[\num{l}][\num{m}]$. از سوی دیگر، بنا بر
\olref[bre]{lem:q-proves-neq}،
$\Th{Q}\Proves\eq/[\num{l}][\num{m}]$. پس~$\Th{Q}$ ناسازگار است.
اما این ناممکن است، زیرا مدل استاندارد~$\Th{Q}$ را ارضا می‌کند
(بنگرید به~\olref[int][def]{def:standard-model})،
$\Sat{N}{\Th{Q}}$؛ و بنا بر قضیهٔ درستی، نظریه‌های ارضاپذیر همواره
سازگارند
(\tagrefs{prfAX/{fol:axd:sou:cor:consistency-soundness},
prfSC/{fol:seq:sou:cor:consistency-soundness},
prfND/{fol:ntd:sou:cor:consistency-soundness},
prfTab/{fol:tab:sou:cor:consistency-soundness}}).
\end{proof}

\begin{lem}
هر تابع بازنمایی‌پذیر در~$\Th{Q}$ محاسبه‌پذیر است.
\end{lem}

\begin{proof}
نخست شهودِ درستی این ادعا را بیان کنیم. برای محاسبهٔ~$f$ چنین
می‌کنیم: همهٔ !!{derivation}sی ممکن~$\delta$ در زبان حساب را فهرست
می‌کنیم؛ این کار را می‌توان به‌طور مکانیکی انجام داد. برای هر یک
می‌آزماییم آیا !!a{derivation} !!a{formula}ی به صورت
$!A_f(\num{n_0},\dots,\num{n_k},\num{m})$ است یا نه (یعنی
!!{formula} بازنمایندهٔ~$f$ در~$\Th{Q}$ از~\olref{lem:rep-q}). اگر چنین
باشد، بنا بر~\olref{lem:rep-q}، $m=f(n_0,\dots,n_k)$ و مقدار~$f$ را
یافته‌ایم. جست‌وجو پایان می‌یابد، زیرا
$\Th{Q}\Proves !A_f(\num{n_0},\dots,\num{n_k},
\num{f(n_0,\dots,n_k)})$؛ پس سرانجام~$\delta$ای از نوع درست می‌یابیم.

این بیان کاملاً دقیق نیست، زیرا روی !!{derivation}s و !!{formula}s عمل
می‌کند، نه فقط روی اعداد؛ و هنوز دقیقاً توضیح نداده‌ایم چرا «فهرست‌کردن
همهٔ !!{derivation}sی ممکن» به‌طور مکانیکی امکان‌پذیر است. اما چنان‌که
دیدیم، می‌توان ترم‌ها، !!{formula}s و !!{derivation}s را با اعداد گودل
رمزگذاری کرد. همچنین مدل دقیقی از محاسبه، یعنی توابع بازگشتی عمومی،
معرفی کرده‌ایم. و دیده‌ایم رابطهٔ~$\Prf[\Th{Q}](d,y)$، که برقرار است
اگر و تنها اگر $d$ عدد گودل !!a{derivation}ی از !!{formula} دارای عدد
گودل~$y$ بر پایهٔ اصول موضوعهٔ~$\Th{Q}$ باشد، (به‌طور اولیه)
بازگشتی است. توابع بازگشتی اولیهٔ دیگری که لازم داریم عبارت‌اند از
$\fn{num}$ (\olref[art][trm]{prop:num-primrec}) و~$\fn{Subst}$
(\olref[art][sub]{prop:subst-primrec}). با استفاده از این‌ها می‌توان
$f$ را با کمینه‌سازی تعریف کرد؛ پس~$f$ بازگشتی است.

نخست تعریف کنید:
\begin{multline*}
  A(n_0, \dots, n_k, m) = \\
  \fn{Subst}(\fn{Subst}(\dots\fn{Subst}(\Gn{!A_f}, \fn{num}(n_0), \Gn{x_0}),\\ \dots),
  \fn{num}(n_k),  \Gn{x_k}), \fn{num}(m), \Gn{y}).
\end{multline*}
این پیچیده به نظر می‌رسد، اما صرفاً تابع زیر است:
$A(n_0,\dots,n_k,m)=
\Gn{!A_f(\num{n_0},\dots,\num{n_k},\num{m})}$.

اکنون رابطهٔ~$R(n_0,\dots,n_k,s)$ را در نظر بگیرید که وقتی~$(s)_0$
عدد گودل !!a{derivation}ی از~$\Th{Q}$ برای
$!A_f(\num{n_0},\dots,\num{n_k},\num{(s)_1})$ باشد برقرار است:
\[
R(n_0, \dots, n_k, s) \quad\text{اگر و تنها اگر}\quad \Prf[\Th{Q}]((s)_0, A(n_0,
\dots, n_k, (s)_1)).
\]
اگر~$s$ای بیابیم که~$R(n_0,\dots,n_k,s)$ درباره‌اش برقرار باشد، زوج
عددی~$(s)_0$ و~$(s)_1$ را یافته‌ایم به‌گونه‌ای که~$(s)_0$ عدد گودل
!!a{derivation}ی از~$A_f(\num{n_0},\dots,\num{n_k},(s)_1)$ است. پس
جست‌وجوی~$s$ مانند جست‌وجوی زوج~$d$ و~$m$ در اثبات غیرصوری است؛ و
تابع محاسبه‌پذیری که چنین~$s$ای را «جست‌وجو می‌کند» با کمینه‌سازی
منظم تعریف‌شدنی است. توجه کنید~$R$ منظم است: به ازای هر
$n_0$،~\dots،~$n_k$، !!a{derivation}ی چون~$\delta$ وجود دارد که گواه
$\Th{Q}\Proves !A_f(\num{n_0},\dots,\num{n_k},
\num{f(n_0,\dots,n_k)})$ است؛ پس $R(n_0,\dots,n_k,s)$ برای
$s=\tuple{\Gn{\delta},f(n_0,\dots,n_k)}$ برقرار است. بنابراین
می‌توانیم~$f$ را چنین بنویسیم:
\[
f(n_0,\dots,n_{k}) = (\umin{s}{R(n_0, \dots, n_k, s)})_1.
\]
\end{proof}

\end{document}
